\section*{Chapter \ref{ch:ml:cross-sectional-deep-models} --- Cross-Sectional Deep Models}

\begin{solution}{exo:ml:cross-sectional-deep-models:1}
$\Gamma$ is $11\times3$: 33 parameters, whatever the number of stocks. Static PCA: $300\times3 = 900$ loadings, and a stock
without history has none.
\end{solution}

\begin{solution}{exo:ml:cross-sectional-deep-models:2}
$(0.5\times2 + (-0.5)(-1) + 1.0\times3)/3 = 1.5\%$. A sum over stocks does not depend on their order.
\end{solution}

\begin{solution}{exo:ml:cross-sectional-deep-models:3}
One-hot into 32 units: $30\times32 = 960$ weights. A one-dimensional embedding: 30 values plus 32 weights from the
embedding to the layer, 62 in all.
\end{solution}

\begin{solution}{exo:ml:cross-sectional-deep-models:4}
PCA's loadings are averages of each stock's betas over the training years, while the betas move with the characteristics:
by the test months the stocks' current betas, and so their expected returns, have little to do with those averages
(correlation 0.03). The total R-squared comes mostly from the first factor, a market factor with betas near one, which any
loading captures.
\end{solution}

\begin{solution}{exo:ml:cross-sectional-deep-models:5}
Predictive R-squared against zero is driven by the level of the forecast relative to the realised mean return; over sixty
months the realised factor means differ from the premia, and a model whose average forecast happens to match them better
scores higher than the truth (1.23\% against 1.09\%). Report total R-squared, the correlation of expected returns with the
truth (on simulated data) or with realised returns (the IC), and their spread over periods.
\end{solution}

\begin{solution}{exo:ml:cross-sectional-deep-models:6}
In-sample fit rises with flexibility and the number of factors; compare the two models out of sample, at the same number
of factors, and with seeds for the \omterm{def:ml:cross-sectional-deep-models:ae}{autoencoder}. Here ten factors against three would also compare different models.
\end{solution}

\begin{solution}{exo:ml:cross-sectional-deep-models:7}
\texttt{embedding(emb\_init=None)}: the embedding's correlation with the premia is 0.05 and the forecasts' correlation with
the truth 0.53 (no better than without the industry). Starting at standard deviation one, the embedding's random initial
values are much larger than the premia's effect on the loss can move them before \omterm{def:ml:trees-and-boosting:early}{early stopping} ends training; started
near zero, what the embedding learns is all it holds.
\end{solution}

\begin{solution}{exo:ml:cross-sectional-deep-models:8}
With $\Gamma = I$ and $K = L$, the factors of month $t$ are $(Z_t^\top Z_t)^{-1}Z_t^\top r_{t+1}$, the least-squares
coefficients of that month's returns on its characteristics: Fama and MacBeth's first step. IPCA with $K < L$ constrains
those coefficients to a $K$-dimensional subspace $\Gamma$.
\end{solution}

\begin{solution}{pb:ml:cross-sectional-deep-models:1}
\begin{enumerate}
\item Total 31.3\%, predictive 1.09\%.
\item Specific returns (8\% a month) are most of the variance, and no factor model explains them.
\item 0.6\%, 0.4\% and 0.3\% a month; they make predictive R-squared mostly a matter of the average level.
\item Half of each beta's cross-sectional variance comes from squares, products and absolute values that no linear function
of the characteristics reproduces.
\item PCA 22.0\%, 22.5\%, 23.0\%; IPCA 25.1\%, 27.6\%, 28.1\%; \omterm{def:ml:cross-sectional-deep-models:ae}{autoencoder} 26.9\%, 30.4\%, 30.8\%.
\item PCA 0.03 and 0.00; IPCA 0.71 and $-0.03$; \omterm{def:ml:cross-sectional-deep-models:ae}{autoencoder} 0.95 and 0.59.
\item Little: 30.4--30.6\% total and 0.95--0.96 correlation across three seeds.
\item Its betas are linear in the characteristics, and the nonlinear part is orthogonal to linear functions of them by
construction.
\item The betas are one function of the characteristics for every stock, so sharing parameters across stocks loses nothing.
\item Test R-squared from 0.004\% to 0.15\%, correlation with the truth from 0.52 to 0.75; the embedding correlates at 0.86 with
the premia against 0.88 for the raw industry averages.
\item Its initialisation: at standard deviation one it never learned the premia.
\item In the factor network's input, the managed portfolios $Z_t^\top r_{t+1}/n$.
\item \emph{Named result}: at three factors, total R-squared 22.5\% (PCA), 27.6\% (IPCA) and 30.4\% (\omterm{def:ml:cross-sectional-deep-models:ae}{autoencoder}) against the
truth's 31.3\%; predictive R-squared 1.01\%, 1.11\% and 1.23\% (truth 1.09\%); the \omterm{def:ml:cross-sectional-deep-models:ae}{autoencoder} recovers 0.59 of the planted
nonlinearity in correlation, IPCA and PCA none.
\item The \omterm{def:ml:cross-sectional-deep-models:ae}{autoencoder}'s betas (or IPCA's when nonlinearity is small): factor exposures must be those of the current
characteristics.
\item On held-out months, by total R-squared, stopping where one more factor no longer helps.
\item Rank-normalise, impute missing values by date, expect the \omterm{def:ml:cross-sectional-deep-models:ae}{autoencoder}'s advantage to shrink with the signal.
\item Nothing from PCA; betas from its characteristics from IPCA and the \omterm{def:ml:cross-sectional-deep-models:ae}{autoencoder}.
\item Seeds, tuning, training time, and interpretability: IPCA's $\Gamma$ can be read.
\item Compare on held-out months at the same number of factors, over several seeds and periods.
\item How a stock's exposures to common risks depend on what the stock is.
\end{enumerate}
\end{solution}

\begin{solution}{iq:ml:cross-sectional-deep-models:1}
PCA gives each stock fixed loadings estimated from its return history; IPCA makes loadings functions of observable
characteristics, so they move as characteristics move and exist for new stocks, with far fewer parameters.
\iqlookfor{time-varying, characteristic-driven betas.}
\end{solution}

\begin{solution}{iq:ml:cross-sectional-deep-models:2}
A beta network maps each stock's characteristics to $K$ betas; a factor network maps the month's characteristic-managed
portfolios to $K$ factors; their product reconstructs the returns, and the loss is the reconstruction error.
\iqlookfor{the two inputs and why the factor input is a portfolio set.}
\end{solution}

\begin{solution}{iq:ml:cross-sectional-deep-models:3}
When stocks have persistent idiosyncratic behaviour not captured by characteristics and enough history each; it can
memorise noise, fails for new identifiers, and leaks through identifiers that change (ticker reuse, mergers). Initialise it
small and regularise.
\iqlookfor{the use case, the failure modes, and practical care.}
\end{solution}

\begin{solution}{iq:ml:cross-sectional-deep-models:4}
Total R-squared compares how well betas explain returns given factors; predictive R-squared over a short period mixes the
cross-section with the level of premia and the period's luck. Use total R-squared and cross-sectional measures (IC, decile
spreads) for comparison.
\iqlookfor{what each measures and its noise.}
\end{solution}

\begin{solution}{iq:ml:cross-sectional-deep-models:5}
Aggregate with a permutation-invariant function: averages of characteristics, characteristic-weighted portfolio returns,
or a learned set encoder $\rho(\sum_i\phi(x_i))$.
\iqlookfor{permutation invariance and fixed-size summaries.}
\end{solution}

\begin{solution}{iq:ml:cross-sectional-deep-models:6}
Given $\Gamma$, each month's factors are the regression of returns on $Z_t\Gamma$; given the factors, $\Gamma$ enters linearly
through $z\otimes f$, and its normal equations are $\sum_t(Z_t^\top Z_t\otimes f_{t+1}f_{t+1}^\top)\operatorname{vec}(\Gamma^\top)
= \sum_tZ_t^\top r_{t+1}\otimes f_{t+1}$; alternate, normalising $\Gamma^\top\Gamma = I$.
\iqlookfor{both steps and the Kronecker structure.}
\end{solution}
